% ============================================================================
% (c) 2025 Mohamed EL AFRIT -- IPSSI
% Licence CC BY-NC-SA 4.0
% https://creativecommons.org/licenses/by-nc-sa/4.0/deed.fr
%
% FICHE DE CORRECTIONS DETAILLEES -- Jour 4 : Reseaux de neurones
% ============================================================================
% ============================================================================
% © 2025 Mohamed EL AFRIT — IPSSI
% Ce contenu est distribué sous licence Creative Commons BY-NC-SA 4.0
% https://creativecommons.org/licenses/by-nc-sa/4.0/deed.fr
% ============================================================================
%
% TEMPLATE COMMUN — Mathématiques appliquées à l'Intelligence Artificielle
% Ce fichier est le préambule partagé de TOUS les documents LaTeX du projet.
%
% Utilisation :
%   % ============================================================================
% © 2025 Mohamed EL AFRIT — IPSSI
% Ce contenu est distribué sous licence Creative Commons BY-NC-SA 4.0
% https://creativecommons.org/licenses/by-nc-sa/4.0/deed.fr
% ============================================================================
%
% TEMPLATE COMMUN — Mathématiques appliquées à l'Intelligence Artificielle
% Ce fichier est le préambule partagé de TOUS les documents LaTeX du projet.
%
% Utilisation :
%   % ============================================================================
% © 2025 Mohamed EL AFRIT — IPSSI
% Ce contenu est distribué sous licence Creative Commons BY-NC-SA 4.0
% https://creativecommons.org/licenses/by-nc-sa/4.0/deed.fr
% ============================================================================
%
% TEMPLATE COMMUN — Mathématiques appliquées à l'Intelligence Artificielle
% Ce fichier est le préambule partagé de TOUS les documents LaTeX du projet.
%
% Utilisation :
%   \input{template_ipssi}          % Importer ce préambule
%   \jourNumero{1}                  % Définir le numéro du jour (1, 2, 3 ou 4)
%   \jourTheme{Données et représentation mathématique}
%   \begin{document}
%   \pagegarde{Fiche de synthèse}{1}{Données et représentation mathématique}
%   ...
%   \end{document}
% ============================================================================

% --- Classe de document ---
\documentclass[11pt,a4paper]{article}

% ============================================================================
% ENCODAGE ET LANGUE
% ============================================================================
\usepackage[utf8]{inputenc}
\usepackage[T1]{fontenc}
\usepackage[french]{babel}

% ============================================================================
% MATHÉMATIQUES
% ============================================================================
\usepackage{amsmath,amssymb,amsfonts,mathtools}

% Commandes mathématiques personnalisées (conventions du cours)
\newcommand{\vx}{\mathbf{x}}          % vecteur x
\newcommand{\vw}{\mathbf{w}}          % vecteur w (poids)
\newcommand{\vy}{\mathbf{y}}          % vecteur y
\newcommand{\mX}{\mathbf{X}}          % matrice X
\newcommand{\mW}{\mathbf{W}}          % matrice W
\newcommand{\yhat}{\hat{y}}           % prédiction
\newcommand{\pscal}[2]{\langle #1, #2 \rangle} % produit scalaire
\newcommand{\gradJ}{\nabla J(\boldsymbol{\theta})} % gradient de J

% ============================================================================
% MISE EN PAGE
% ============================================================================
\usepackage[margin=2.5cm]{geometry}
\usepackage{setspace}
\setstretch{1.15}                      % interligne légèrement aéré

% ============================================================================
% COULEURS
% ============================================================================
\usepackage[dvipsnames,svgnames,x11names]{xcolor}

% Palette principale du projet
\definecolor{ipssiBleu}{HTML}{3B82F6}       % Définitions
\definecolor{ipssiVert}{HTML}{22C55E}       % À retenir
\definecolor{ipssiOrange}{HTML}{F97316}     % Attention / Piège
\definecolor{ipssiViolet}{HTML}{A855F7}     % Intuition
\definecolor{ipssiCyan}{HTML}{06B6D4}       % Exemples
\definecolor{ipssiGrisClair}{HTML}{F3F4F6}  % Fond code
\definecolor{ipssiGrisFonce}{HTML}{1F2937}  % Texte code
\definecolor{ipssiFondPage}{HTML}{FAFAFA}   % Fond léger page de garde
\definecolor{ipssiRouge}{HTML}{E74C3C}      % Classe 0 (salaire ≤ 45k€)
\definecolor{ipssiVertData}{HTML}{2ECC71}   % Classe 1 (salaire > 45k€)
\definecolor{ipssiBleuDroite}{HTML}{3498DB} % Droite de régression

% ============================================================================
% ENCADRÉS PÉDAGOGIQUES (tcolorbox)
% ============================================================================
\usepackage[most]{tcolorbox}

% Style de base partagé
\tcbset{
    ipssibase/.style={
        enhanced,
        breakable,
        left=4mm,
        right=4mm,
        top=3mm,
        bottom=3mm,
        boxrule=0pt,
        frame hidden,
        borderline west={3pt}{0pt}{#1},
        colback=#1!5,
        colframe=#1,
        fonttitle=\bfseries\sffamily,
        coltitle=#1!80!black,
        attach boxed title to top left={yshift=-2mm, xshift=4mm},
        boxed title style={
            colback=#1!15,
            colframe=#1,
            boxrule=0.5pt,
            arc=2pt,
        },
        arc=2pt,
        outer arc=2pt,
    }
}

% Icônes TikZ pour les encadrés (compatible pdflatex, pas d'emoji Unicode)
\newcommand{\icoDefinition}{%
    \tikz[baseline=-0.3em]{\draw[ipssiBleu, thick] (0,0) -- (0.3,0) -- (0.3,0.35) -- (0,0.35) -- cycle;
    \draw[ipssiBleu, thick] (0.05,0.1) -- (0.25,0.1); \draw[ipssiBleu, thick] (0.05,0.2) -- (0.2,0.2);}%
}
\newcommand{\icoRetenir}{%
    \tikz[baseline=-0.2em]{\draw[ipssiVert, thick, fill=ipssiVert!20] (0.15,0) -- (0.3,0.3) -- (0,0.3) -- cycle;}%
}
\newcommand{\icoAttention}{%
    \tikz[baseline=-0.2em]{\draw[ipssiOrange, thick] (0.15,0.35) -- (0,0) -- (0.3,0) -- cycle;
    \node[ipssiOrange, font=\tiny\bfseries] at (0.15,0.12) {!};}%
}
\newcommand{\icoIntuition}{%
    \tikz[baseline=-0.2em]{\draw[ipssiViolet, thick, fill=ipssiViolet!15]
    (0.15,0.35) -- (0.08,0.2) -- (0.08,0.15) -- (0.22,0.15) -- (0.22,0.2) -- cycle;
    \draw[ipssiViolet, thick] (0.11,0.12) -- (0.11,0.05); \draw[ipssiViolet, thick] (0.19,0.12) -- (0.19,0.05);}%
}
\newcommand{\icoExemple}{%
    \tikz[baseline=-0.2em]{\draw[ipssiCyan, thick] (0,0) rectangle (0.3,0.3);
    \draw[ipssiCyan, thick] (0,0.15) -- (0.3,0.15); \draw[ipssiCyan, thick] (0.15,0) -- (0.15,0.3);}%
}
\newcommand{\icoPython}{%
    \tikz[baseline=-0.2em]{\node[font=\small\bfseries\ttfamily, ipssiVert!70!black] at (0,0) {\textgreater\textgreater};}%
}

% Définition (bleu)
\newtcolorbox{definitionbox}[1][D\'{e}finition]{
    ipssibase=ipssiBleu,
    title={\icoDefinition\; #1},
}

% À retenir (vert)
\newtcolorbox{retenirbox}[1][\`{A} retenir]{
    ipssibase=ipssiVert,
    title={\icoRetenir\; #1},
}

% Attention / Piège (orange)
\newtcolorbox{attentionbox}[1][Attention]{
    ipssibase=ipssiOrange,
    title={\icoAttention\; #1},
}

% Intuition (violet)
\newtcolorbox{intuitionbox}[1][Intuition]{
    ipssibase=ipssiViolet,
    title={\icoIntuition\; #1},
}

% Exemple (cyan)
\newtcolorbox{exemplebox}[1][Exemple]{
    ipssibase=ipssiCyan,
    title={\icoExemple\; #1},
}

% Code Python (style terminal)
\newtcolorbox{pythonbox}[1][Code Python]{
    enhanced,
    breakable,
    left=4mm, right=4mm, top=3mm, bottom=3mm,
    boxrule=0.5pt,
    colback=ipssiGrisClair,
    colframe=ipssiGrisFonce!40,
    fonttitle=\bfseries\ttfamily\small,
    coltitle=ipssiGrisFonce,
    title={\icoPython\; #1},
    attach boxed title to top left={yshift=-2mm, xshift=4mm},
    boxed title style={colback=ipssiGrisClair, colframe=ipssiGrisFonce!40, boxrule=0.5pt, arc=2pt},
    arc=2pt, outer arc=2pt,
}

% ============================================================================
% CODE PYTHON (listings)
% ============================================================================
\usepackage{listings}

\lstdefinestyle{python}{
    language=Python,
    basicstyle=\ttfamily\small,
    keywordstyle=\color{ipssiBleu}\bfseries,
    stringstyle=\color{ipssiVert!70!black},
    commentstyle=\color{gray}\itshape,
    numberstyle=\tiny\color{gray},
    numbers=left,
    numbersep=8pt,
    backgroundcolor=\color{ipssiGrisClair},
    frame=none,
    rulecolor=\color{ipssiGrisFonce!20},
    breaklines=true,
    breakatwhitespace=true,
    tabsize=4,
    showstringspaces=false,
    captionpos=b,
    aboveskip=8pt,
    belowskip=8pt,
    xleftmargin=10pt,
    framexleftmargin=10pt,
    morekeywords={as, with, True, False, None, self, np, plt, import, from},
    literate=
        {é}{{\'e}}1 {è}{{\`e}}1 {ê}{{\^e}}1 {ë}{{\"e}}1
        {à}{{\`a}}1 {â}{{\^a}}1
        {ù}{{\`u}}1 {û}{{\^u}}1
        {ô}{{\^o}}1
        {î}{{\^i}}1 {ï}{{\"i}}1
        {ç}{{\c{c}}}1
        {É}{{\'E}}1 {È}{{\`E}}1
        {→}{{$\rightarrow$}}1
        {≈}{{$\approx$}}1
        {≤}{{$\leq$}}1
        {≥}{{$\geq$}}1
        {★}{{$\bigstar$}}1
        {ŷ}{{$\hat{y}$}}1
        {·}{{$\cdot$}}1,
}

\lstset{style=python}

% Style pour le code inline
\newcommand{\pyinline}[1]{\lstinline[style=python,basicstyle=\ttfamily\small]{#1}}

% ============================================================================
% GRAPHIQUES
% ============================================================================
\usepackage{tikz}
\usepackage{pgfplots}
\pgfplotsset{compat=1.18}
\usetikzlibrary{arrows.meta, calc, positioning, shapes.geometric, decorations.pathreplacing}

% ============================================================================
% TABLEAUX
% ============================================================================
\usepackage{booktabs}
\usepackage{tabularx}
\usepackage{multirow}

% ============================================================================
% LISTES
% ============================================================================
\usepackage{enumitem}
\setlist[itemize]{itemsep=2pt, parsep=0pt, topsep=4pt}
\setlist[enumerate]{itemsep=2pt, parsep=0pt, topsep=4pt}

% ============================================================================
% HYPERLIENS
% ============================================================================
\usepackage[
    colorlinks=true,
    linkcolor=ipssiBleu!80!black,
    urlcolor=ipssiBleu!80!black,
    citecolor=ipssiVert!80!black,
    bookmarks=true,
    bookmarksopen=true,
    pdfauthor={Mohamed EL AFRIT},
    pdftitle={Mathématiques appliquées à l'IA — IPSSI},
    pdfsubject={Formation Bachelor 2 — IPSSI 2025-2026},
]{hyperref}

% ============================================================================
% IMAGES
% ============================================================================
\usepackage{graphicx}

% ============================================================================
% VARIABLES PARAMÉTRABLES (jour et thème)
% ============================================================================
\newcommand{\leJour}{X}
\newcommand{\leTheme}{Thème}
\newcommand{\jourNumero}[1]{\renewcommand{\leJour}{#1}}
\newcommand{\jourTheme}[1]{\renewcommand{\leTheme}{#1}}

% ============================================================================
% EN-TÊTE ET PIED DE PAGE (fancyhdr)
% ============================================================================
\usepackage{fancyhdr}

\pagestyle{fancy}
\fancyhf{}

% En-tête
\fancyhead[L]{\small\sffamily\textcolor{gray}{IPSSI — Maths appliquées à l'IA}}
\fancyhead[R]{\small\sffamily\textcolor{gray}{Jour \leJour\ — \leTheme}}
\renewcommand{\headrulewidth}{0.4pt}

% Pied de page
\fancyfoot[L]{\footnotesize\textcolor{gray}{Mohamed EL AFRIT\enspace\textbullet\enspace\url{www.mohamedelafrit.com}\enspace\textbullet\enspace CC BY-NC-SA 4.0}}
\fancyfoot[C]{\footnotesize\textcolor{gray}{\thepage}}
\fancyfoot[R]{\footnotesize\textcolor{gray}{2025–2026}}
\renewcommand{\footrulewidth}{0.2pt}

% Style pour la page de garde (pas d'en-tête)
\fancypagestyle{pagegarde}{
    \fancyhf{}
    \renewcommand{\headrulewidth}{0pt}
    \fancyfoot[C]{\footnotesize\textcolor{gray}{\thepage}}
    \renewcommand{\footrulewidth}{0pt}
}

% ============================================================================
% PAGE DE GARDE — \pagegarde{titre}{jour}{theme}
% ============================================================================
\newcommand{\pagegarde}[3]{%
    \thispagestyle{pagegarde}
    \begin{center}
        \vspace*{1.5cm}

        % Logo / Bandeau institution
        {\Large\sffamily\bfseries\textcolor{ipssiBleu}{IPSSI}}\par
        \vspace{2pt}
        {\normalsize\sffamily 2\textsuperscript{e} année Bachelor Informatique}\par

        \vspace{0.8cm}
        {\color{ipssiBleu}\rule{0.6\textwidth}{1.5pt}}\par
        \vspace{0.8cm}

        % Titre du module
        {\LARGE\sffamily\bfseries Mathématiques appliquées\\[4pt] à l'Intelligence Artificielle}\par

        \vspace{1cm}

        % Titre du document
        \begin{tcolorbox}[
            enhanced,
            colback=ipssiBleu!5,
            colframe=ipssiBleu!50,
            boxrule=1pt,
            arc=4pt,
            width=0.85\textwidth,
            halign=center,
            top=10pt, bottom=10pt,
        ]
            {\huge\sffamily\bfseries\textcolor{ipssiBleu!80!black}{#1}}\par
            \vspace{6pt}
            {\Large\sffamily Jour #2 — #3}
        \end{tcolorbox}

        \vspace{1.5cm}

        % Informations auteur
        {\large\sffamily
            \textbf{Auteur :} Mohamed EL AFRIT\\[4pt]
            \url{www.mohamedelafrit.com}\\[8pt]
            \textbf{Année académique :} 2025–2026
        }\par

        \vspace{1.5cm}

        % Licence Creative Commons
        {\color{ipssiBleu}\rule{0.4\textwidth}{0.5pt}}\par
        \vspace{8pt}
        {\small\sffamily
            \textcopyright\ 2025 Mohamed EL AFRIT — IPSSI\\[2pt]
            Ce contenu est distribué sous licence\\[2pt]
            \textbf{Creative Commons BY-NC-SA 4.0}\\[2pt]
            \url{https://creativecommons.org/licenses/by-nc-sa/4.0/deed.fr}
        }\par
    \end{center}
    \newpage
}

% ============================================================================
% NIVEAUX D'EXERCICES
% ============================================================================

% Étoile compatible pdflatex
\newcommand{\starfull}[1]{\textcolor{#1}{$\bigstar$}}
\newcommand{\starempty}{\textcolor{gray!40}{$\bigstar$}}

% Fondamental (1 étoile)
\newcommand{\exofondamental}{%
    \starfull{ipssiVert}\starempty\starempty\enspace%
    {\small\sffamily\textcolor{ipssiVert}{\textbf{Fondamental}}}%
}

% Intermédiaire (2 étoiles)
\newcommand{\exointermediaire}{%
    \starfull{ipssiOrange}\starfull{ipssiOrange}\starempty\enspace%
    {\small\sffamily\textcolor{ipssiOrange}{\textbf{Interm\'{e}diaire}}}%
}

% Avancé (3 étoiles)
\newcommand{\exoavance}{%
    \starfull{ipssiRouge}\starfull{ipssiRouge}\starfull{ipssiRouge}\enspace%
    {\small\sffamily\textcolor{ipssiRouge}{\textbf{Avanc\'{e}}}}%
}

% Commande d'exercice numéroté
\newcounter{exocount}[section]
\newcommand{\exercice}[1]{%
    \stepcounter{exocount}%
    \vspace{8pt}%
    \noindent\textbf{\sffamily Exercice \theexocount}\enspace — \enspace #1\par
    \vspace{4pt}%
}

% ============================================================================
% COMMANDES UTILITAIRES
% ============================================================================

% Séparateur visuel entre sections
\newcommand{\separateur}{%
    \vspace{6pt}
    \begin{center}
        {\color{ipssiBleu!30}\rule{0.3\textwidth}{0.5pt}}
    \end{center}
    \vspace{6pt}
}

% Mot-clé en gras coloré
\newcommand{\motcle}[1]{\textbf{\textcolor{ipssiBleu!80!black}{#1}}}

% Formule encadrée importante
\newcommand{\formule}[1]{%
    \begin{center}
        \fcolorbox{ipssiBleu!50}{ipssiBleu!5}{%
            \quad$\displaystyle #1$\quad%
        }
    \end{center}
}

% ============================================================================
% FIN DU PRÉAMBULE
% ============================================================================
          % Importer ce préambule
%   \jourNumero{1}                  % Définir le numéro du jour (1, 2, 3 ou 4)
%   \jourTheme{Données et représentation mathématique}
%   \begin{document}
%   \pagegarde{Fiche de synthèse}{1}{Données et représentation mathématique}
%   ...
%   \end{document}
% ============================================================================

% --- Classe de document ---
\documentclass[11pt,a4paper]{article}

% ============================================================================
% ENCODAGE ET LANGUE
% ============================================================================
\usepackage[utf8]{inputenc}
\usepackage[T1]{fontenc}
\usepackage[french]{babel}

% ============================================================================
% MATHÉMATIQUES
% ============================================================================
\usepackage{amsmath,amssymb,amsfonts,mathtools}

% Commandes mathématiques personnalisées (conventions du cours)
\newcommand{\vx}{\mathbf{x}}          % vecteur x
\newcommand{\vw}{\mathbf{w}}          % vecteur w (poids)
\newcommand{\vy}{\mathbf{y}}          % vecteur y
\newcommand{\mX}{\mathbf{X}}          % matrice X
\newcommand{\mW}{\mathbf{W}}          % matrice W
\newcommand{\yhat}{\hat{y}}           % prédiction
\newcommand{\pscal}[2]{\langle #1, #2 \rangle} % produit scalaire
\newcommand{\gradJ}{\nabla J(\boldsymbol{\theta})} % gradient de J

% ============================================================================
% MISE EN PAGE
% ============================================================================
\usepackage[margin=2.5cm]{geometry}
\usepackage{setspace}
\setstretch{1.15}                      % interligne légèrement aéré

% ============================================================================
% COULEURS
% ============================================================================
\usepackage[dvipsnames,svgnames,x11names]{xcolor}

% Palette principale du projet
\definecolor{ipssiBleu}{HTML}{3B82F6}       % Définitions
\definecolor{ipssiVert}{HTML}{22C55E}       % À retenir
\definecolor{ipssiOrange}{HTML}{F97316}     % Attention / Piège
\definecolor{ipssiViolet}{HTML}{A855F7}     % Intuition
\definecolor{ipssiCyan}{HTML}{06B6D4}       % Exemples
\definecolor{ipssiGrisClair}{HTML}{F3F4F6}  % Fond code
\definecolor{ipssiGrisFonce}{HTML}{1F2937}  % Texte code
\definecolor{ipssiFondPage}{HTML}{FAFAFA}   % Fond léger page de garde
\definecolor{ipssiRouge}{HTML}{E74C3C}      % Classe 0 (salaire ≤ 45k€)
\definecolor{ipssiVertData}{HTML}{2ECC71}   % Classe 1 (salaire > 45k€)
\definecolor{ipssiBleuDroite}{HTML}{3498DB} % Droite de régression

% ============================================================================
% ENCADRÉS PÉDAGOGIQUES (tcolorbox)
% ============================================================================
\usepackage[most]{tcolorbox}

% Style de base partagé
\tcbset{
    ipssibase/.style={
        enhanced,
        breakable,
        left=4mm,
        right=4mm,
        top=3mm,
        bottom=3mm,
        boxrule=0pt,
        frame hidden,
        borderline west={3pt}{0pt}{#1},
        colback=#1!5,
        colframe=#1,
        fonttitle=\bfseries\sffamily,
        coltitle=#1!80!black,
        attach boxed title to top left={yshift=-2mm, xshift=4mm},
        boxed title style={
            colback=#1!15,
            colframe=#1,
            boxrule=0.5pt,
            arc=2pt,
        },
        arc=2pt,
        outer arc=2pt,
    }
}

% Icônes TikZ pour les encadrés (compatible pdflatex, pas d'emoji Unicode)
\newcommand{\icoDefinition}{%
    \tikz[baseline=-0.3em]{\draw[ipssiBleu, thick] (0,0) -- (0.3,0) -- (0.3,0.35) -- (0,0.35) -- cycle;
    \draw[ipssiBleu, thick] (0.05,0.1) -- (0.25,0.1); \draw[ipssiBleu, thick] (0.05,0.2) -- (0.2,0.2);}%
}
\newcommand{\icoRetenir}{%
    \tikz[baseline=-0.2em]{\draw[ipssiVert, thick, fill=ipssiVert!20] (0.15,0) -- (0.3,0.3) -- (0,0.3) -- cycle;}%
}
\newcommand{\icoAttention}{%
    \tikz[baseline=-0.2em]{\draw[ipssiOrange, thick] (0.15,0.35) -- (0,0) -- (0.3,0) -- cycle;
    \node[ipssiOrange, font=\tiny\bfseries] at (0.15,0.12) {!};}%
}
\newcommand{\icoIntuition}{%
    \tikz[baseline=-0.2em]{\draw[ipssiViolet, thick, fill=ipssiViolet!15]
    (0.15,0.35) -- (0.08,0.2) -- (0.08,0.15) -- (0.22,0.15) -- (0.22,0.2) -- cycle;
    \draw[ipssiViolet, thick] (0.11,0.12) -- (0.11,0.05); \draw[ipssiViolet, thick] (0.19,0.12) -- (0.19,0.05);}%
}
\newcommand{\icoExemple}{%
    \tikz[baseline=-0.2em]{\draw[ipssiCyan, thick] (0,0) rectangle (0.3,0.3);
    \draw[ipssiCyan, thick] (0,0.15) -- (0.3,0.15); \draw[ipssiCyan, thick] (0.15,0) -- (0.15,0.3);}%
}
\newcommand{\icoPython}{%
    \tikz[baseline=-0.2em]{\node[font=\small\bfseries\ttfamily, ipssiVert!70!black] at (0,0) {\textgreater\textgreater};}%
}

% Définition (bleu)
\newtcolorbox{definitionbox}[1][D\'{e}finition]{
    ipssibase=ipssiBleu,
    title={\icoDefinition\; #1},
}

% À retenir (vert)
\newtcolorbox{retenirbox}[1][\`{A} retenir]{
    ipssibase=ipssiVert,
    title={\icoRetenir\; #1},
}

% Attention / Piège (orange)
\newtcolorbox{attentionbox}[1][Attention]{
    ipssibase=ipssiOrange,
    title={\icoAttention\; #1},
}

% Intuition (violet)
\newtcolorbox{intuitionbox}[1][Intuition]{
    ipssibase=ipssiViolet,
    title={\icoIntuition\; #1},
}

% Exemple (cyan)
\newtcolorbox{exemplebox}[1][Exemple]{
    ipssibase=ipssiCyan,
    title={\icoExemple\; #1},
}

% Code Python (style terminal)
\newtcolorbox{pythonbox}[1][Code Python]{
    enhanced,
    breakable,
    left=4mm, right=4mm, top=3mm, bottom=3mm,
    boxrule=0.5pt,
    colback=ipssiGrisClair,
    colframe=ipssiGrisFonce!40,
    fonttitle=\bfseries\ttfamily\small,
    coltitle=ipssiGrisFonce,
    title={\icoPython\; #1},
    attach boxed title to top left={yshift=-2mm, xshift=4mm},
    boxed title style={colback=ipssiGrisClair, colframe=ipssiGrisFonce!40, boxrule=0.5pt, arc=2pt},
    arc=2pt, outer arc=2pt,
}

% ============================================================================
% CODE PYTHON (listings)
% ============================================================================
\usepackage{listings}

\lstdefinestyle{python}{
    language=Python,
    basicstyle=\ttfamily\small,
    keywordstyle=\color{ipssiBleu}\bfseries,
    stringstyle=\color{ipssiVert!70!black},
    commentstyle=\color{gray}\itshape,
    numberstyle=\tiny\color{gray},
    numbers=left,
    numbersep=8pt,
    backgroundcolor=\color{ipssiGrisClair},
    frame=none,
    rulecolor=\color{ipssiGrisFonce!20},
    breaklines=true,
    breakatwhitespace=true,
    tabsize=4,
    showstringspaces=false,
    captionpos=b,
    aboveskip=8pt,
    belowskip=8pt,
    xleftmargin=10pt,
    framexleftmargin=10pt,
    morekeywords={as, with, True, False, None, self, np, plt, import, from},
    literate=
        {é}{{\'e}}1 {è}{{\`e}}1 {ê}{{\^e}}1 {ë}{{\"e}}1
        {à}{{\`a}}1 {â}{{\^a}}1
        {ù}{{\`u}}1 {û}{{\^u}}1
        {ô}{{\^o}}1
        {î}{{\^i}}1 {ï}{{\"i}}1
        {ç}{{\c{c}}}1
        {É}{{\'E}}1 {È}{{\`E}}1
        {→}{{$\rightarrow$}}1
        {≈}{{$\approx$}}1
        {≤}{{$\leq$}}1
        {≥}{{$\geq$}}1
        {★}{{$\bigstar$}}1
        {ŷ}{{$\hat{y}$}}1
        {·}{{$\cdot$}}1,
}

\lstset{style=python}

% Style pour le code inline
\newcommand{\pyinline}[1]{\lstinline[style=python,basicstyle=\ttfamily\small]{#1}}

% ============================================================================
% GRAPHIQUES
% ============================================================================
\usepackage{tikz}
\usepackage{pgfplots}
\pgfplotsset{compat=1.18}
\usetikzlibrary{arrows.meta, calc, positioning, shapes.geometric, decorations.pathreplacing}

% ============================================================================
% TABLEAUX
% ============================================================================
\usepackage{booktabs}
\usepackage{tabularx}
\usepackage{multirow}

% ============================================================================
% LISTES
% ============================================================================
\usepackage{enumitem}
\setlist[itemize]{itemsep=2pt, parsep=0pt, topsep=4pt}
\setlist[enumerate]{itemsep=2pt, parsep=0pt, topsep=4pt}

% ============================================================================
% HYPERLIENS
% ============================================================================
\usepackage[
    colorlinks=true,
    linkcolor=ipssiBleu!80!black,
    urlcolor=ipssiBleu!80!black,
    citecolor=ipssiVert!80!black,
    bookmarks=true,
    bookmarksopen=true,
    pdfauthor={Mohamed EL AFRIT},
    pdftitle={Mathématiques appliquées à l'IA — IPSSI},
    pdfsubject={Formation Bachelor 2 — IPSSI 2025-2026},
]{hyperref}

% ============================================================================
% IMAGES
% ============================================================================
\usepackage{graphicx}

% ============================================================================
% VARIABLES PARAMÉTRABLES (jour et thème)
% ============================================================================
\newcommand{\leJour}{X}
\newcommand{\leTheme}{Thème}
\newcommand{\jourNumero}[1]{\renewcommand{\leJour}{#1}}
\newcommand{\jourTheme}[1]{\renewcommand{\leTheme}{#1}}

% ============================================================================
% EN-TÊTE ET PIED DE PAGE (fancyhdr)
% ============================================================================
\usepackage{fancyhdr}

\pagestyle{fancy}
\fancyhf{}

% En-tête
\fancyhead[L]{\small\sffamily\textcolor{gray}{IPSSI — Maths appliquées à l'IA}}
\fancyhead[R]{\small\sffamily\textcolor{gray}{Jour \leJour\ — \leTheme}}
\renewcommand{\headrulewidth}{0.4pt}

% Pied de page
\fancyfoot[L]{\footnotesize\textcolor{gray}{Mohamed EL AFRIT\enspace\textbullet\enspace\url{www.mohamedelafrit.com}\enspace\textbullet\enspace CC BY-NC-SA 4.0}}
\fancyfoot[C]{\footnotesize\textcolor{gray}{\thepage}}
\fancyfoot[R]{\footnotesize\textcolor{gray}{2025–2026}}
\renewcommand{\footrulewidth}{0.2pt}

% Style pour la page de garde (pas d'en-tête)
\fancypagestyle{pagegarde}{
    \fancyhf{}
    \renewcommand{\headrulewidth}{0pt}
    \fancyfoot[C]{\footnotesize\textcolor{gray}{\thepage}}
    \renewcommand{\footrulewidth}{0pt}
}

% ============================================================================
% PAGE DE GARDE — \pagegarde{titre}{jour}{theme}
% ============================================================================
\newcommand{\pagegarde}[3]{%
    \thispagestyle{pagegarde}
    \begin{center}
        \vspace*{1.5cm}

        % Logo / Bandeau institution
        {\Large\sffamily\bfseries\textcolor{ipssiBleu}{IPSSI}}\par
        \vspace{2pt}
        {\normalsize\sffamily 2\textsuperscript{e} année Bachelor Informatique}\par

        \vspace{0.8cm}
        {\color{ipssiBleu}\rule{0.6\textwidth}{1.5pt}}\par
        \vspace{0.8cm}

        % Titre du module
        {\LARGE\sffamily\bfseries Mathématiques appliquées\\[4pt] à l'Intelligence Artificielle}\par

        \vspace{1cm}

        % Titre du document
        \begin{tcolorbox}[
            enhanced,
            colback=ipssiBleu!5,
            colframe=ipssiBleu!50,
            boxrule=1pt,
            arc=4pt,
            width=0.85\textwidth,
            halign=center,
            top=10pt, bottom=10pt,
        ]
            {\huge\sffamily\bfseries\textcolor{ipssiBleu!80!black}{#1}}\par
            \vspace{6pt}
            {\Large\sffamily Jour #2 — #3}
        \end{tcolorbox}

        \vspace{1.5cm}

        % Informations auteur
        {\large\sffamily
            \textbf{Auteur :} Mohamed EL AFRIT\\[4pt]
            \url{www.mohamedelafrit.com}\\[8pt]
            \textbf{Année académique :} 2025–2026
        }\par

        \vspace{1.5cm}

        % Licence Creative Commons
        {\color{ipssiBleu}\rule{0.4\textwidth}{0.5pt}}\par
        \vspace{8pt}
        {\small\sffamily
            \textcopyright\ 2025 Mohamed EL AFRIT — IPSSI\\[2pt]
            Ce contenu est distribué sous licence\\[2pt]
            \textbf{Creative Commons BY-NC-SA 4.0}\\[2pt]
            \url{https://creativecommons.org/licenses/by-nc-sa/4.0/deed.fr}
        }\par
    \end{center}
    \newpage
}

% ============================================================================
% NIVEAUX D'EXERCICES
% ============================================================================

% Étoile compatible pdflatex
\newcommand{\starfull}[1]{\textcolor{#1}{$\bigstar$}}
\newcommand{\starempty}{\textcolor{gray!40}{$\bigstar$}}

% Fondamental (1 étoile)
\newcommand{\exofondamental}{%
    \starfull{ipssiVert}\starempty\starempty\enspace%
    {\small\sffamily\textcolor{ipssiVert}{\textbf{Fondamental}}}%
}

% Intermédiaire (2 étoiles)
\newcommand{\exointermediaire}{%
    \starfull{ipssiOrange}\starfull{ipssiOrange}\starempty\enspace%
    {\small\sffamily\textcolor{ipssiOrange}{\textbf{Interm\'{e}diaire}}}%
}

% Avancé (3 étoiles)
\newcommand{\exoavance}{%
    \starfull{ipssiRouge}\starfull{ipssiRouge}\starfull{ipssiRouge}\enspace%
    {\small\sffamily\textcolor{ipssiRouge}{\textbf{Avanc\'{e}}}}%
}

% Commande d'exercice numéroté
\newcounter{exocount}[section]
\newcommand{\exercice}[1]{%
    \stepcounter{exocount}%
    \vspace{8pt}%
    \noindent\textbf{\sffamily Exercice \theexocount}\enspace — \enspace #1\par
    \vspace{4pt}%
}

% ============================================================================
% COMMANDES UTILITAIRES
% ============================================================================

% Séparateur visuel entre sections
\newcommand{\separateur}{%
    \vspace{6pt}
    \begin{center}
        {\color{ipssiBleu!30}\rule{0.3\textwidth}{0.5pt}}
    \end{center}
    \vspace{6pt}
}

% Mot-clé en gras coloré
\newcommand{\motcle}[1]{\textbf{\textcolor{ipssiBleu!80!black}{#1}}}

% Formule encadrée importante
\newcommand{\formule}[1]{%
    \begin{center}
        \fcolorbox{ipssiBleu!50}{ipssiBleu!5}{%
            \quad$\displaystyle #1$\quad%
        }
    \end{center}
}

% ============================================================================
% FIN DU PRÉAMBULE
% ============================================================================
          % Importer ce préambule
%   \jourNumero{1}                  % Définir le numéro du jour (1, 2, 3 ou 4)
%   \jourTheme{Données et représentation mathématique}
%   \begin{document}
%   \pagegarde{Fiche de synthèse}{1}{Données et représentation mathématique}
%   ...
%   \end{document}
% ============================================================================

% --- Classe de document ---
\documentclass[11pt,a4paper]{article}

% ============================================================================
% ENCODAGE ET LANGUE
% ============================================================================
\usepackage[utf8]{inputenc}
\usepackage[T1]{fontenc}
\usepackage[french]{babel}

% ============================================================================
% MATHÉMATIQUES
% ============================================================================
\usepackage{amsmath,amssymb,amsfonts,mathtools}

% Commandes mathématiques personnalisées (conventions du cours)
\newcommand{\vx}{\mathbf{x}}          % vecteur x
\newcommand{\vw}{\mathbf{w}}          % vecteur w (poids)
\newcommand{\vy}{\mathbf{y}}          % vecteur y
\newcommand{\mX}{\mathbf{X}}          % matrice X
\newcommand{\mW}{\mathbf{W}}          % matrice W
\newcommand{\yhat}{\hat{y}}           % prédiction
\newcommand{\pscal}[2]{\langle #1, #2 \rangle} % produit scalaire
\newcommand{\gradJ}{\nabla J(\boldsymbol{\theta})} % gradient de J

% ============================================================================
% MISE EN PAGE
% ============================================================================
\usepackage[margin=2.5cm]{geometry}
\usepackage{setspace}
\setstretch{1.15}                      % interligne légèrement aéré

% ============================================================================
% COULEURS
% ============================================================================
\usepackage[dvipsnames,svgnames,x11names]{xcolor}

% Palette principale du projet
\definecolor{ipssiBleu}{HTML}{3B82F6}       % Définitions
\definecolor{ipssiVert}{HTML}{22C55E}       % À retenir
\definecolor{ipssiOrange}{HTML}{F97316}     % Attention / Piège
\definecolor{ipssiViolet}{HTML}{A855F7}     % Intuition
\definecolor{ipssiCyan}{HTML}{06B6D4}       % Exemples
\definecolor{ipssiGrisClair}{HTML}{F3F4F6}  % Fond code
\definecolor{ipssiGrisFonce}{HTML}{1F2937}  % Texte code
\definecolor{ipssiFondPage}{HTML}{FAFAFA}   % Fond léger page de garde
\definecolor{ipssiRouge}{HTML}{E74C3C}      % Classe 0 (salaire ≤ 45k€)
\definecolor{ipssiVertData}{HTML}{2ECC71}   % Classe 1 (salaire > 45k€)
\definecolor{ipssiBleuDroite}{HTML}{3498DB} % Droite de régression

% ============================================================================
% ENCADRÉS PÉDAGOGIQUES (tcolorbox)
% ============================================================================
\usepackage[most]{tcolorbox}

% Style de base partagé
\tcbset{
    ipssibase/.style={
        enhanced,
        breakable,
        left=4mm,
        right=4mm,
        top=3mm,
        bottom=3mm,
        boxrule=0pt,
        frame hidden,
        borderline west={3pt}{0pt}{#1},
        colback=#1!5,
        colframe=#1,
        fonttitle=\bfseries\sffamily,
        coltitle=#1!80!black,
        attach boxed title to top left={yshift=-2mm, xshift=4mm},
        boxed title style={
            colback=#1!15,
            colframe=#1,
            boxrule=0.5pt,
            arc=2pt,
        },
        arc=2pt,
        outer arc=2pt,
    }
}

% Icônes TikZ pour les encadrés (compatible pdflatex, pas d'emoji Unicode)
\newcommand{\icoDefinition}{%
    \tikz[baseline=-0.3em]{\draw[ipssiBleu, thick] (0,0) -- (0.3,0) -- (0.3,0.35) -- (0,0.35) -- cycle;
    \draw[ipssiBleu, thick] (0.05,0.1) -- (0.25,0.1); \draw[ipssiBleu, thick] (0.05,0.2) -- (0.2,0.2);}%
}
\newcommand{\icoRetenir}{%
    \tikz[baseline=-0.2em]{\draw[ipssiVert, thick, fill=ipssiVert!20] (0.15,0) -- (0.3,0.3) -- (0,0.3) -- cycle;}%
}
\newcommand{\icoAttention}{%
    \tikz[baseline=-0.2em]{\draw[ipssiOrange, thick] (0.15,0.35) -- (0,0) -- (0.3,0) -- cycle;
    \node[ipssiOrange, font=\tiny\bfseries] at (0.15,0.12) {!};}%
}
\newcommand{\icoIntuition}{%
    \tikz[baseline=-0.2em]{\draw[ipssiViolet, thick, fill=ipssiViolet!15]
    (0.15,0.35) -- (0.08,0.2) -- (0.08,0.15) -- (0.22,0.15) -- (0.22,0.2) -- cycle;
    \draw[ipssiViolet, thick] (0.11,0.12) -- (0.11,0.05); \draw[ipssiViolet, thick] (0.19,0.12) -- (0.19,0.05);}%
}
\newcommand{\icoExemple}{%
    \tikz[baseline=-0.2em]{\draw[ipssiCyan, thick] (0,0) rectangle (0.3,0.3);
    \draw[ipssiCyan, thick] (0,0.15) -- (0.3,0.15); \draw[ipssiCyan, thick] (0.15,0) -- (0.15,0.3);}%
}
\newcommand{\icoPython}{%
    \tikz[baseline=-0.2em]{\node[font=\small\bfseries\ttfamily, ipssiVert!70!black] at (0,0) {\textgreater\textgreater};}%
}

% Définition (bleu)
\newtcolorbox{definitionbox}[1][D\'{e}finition]{
    ipssibase=ipssiBleu,
    title={\icoDefinition\; #1},
}

% À retenir (vert)
\newtcolorbox{retenirbox}[1][\`{A} retenir]{
    ipssibase=ipssiVert,
    title={\icoRetenir\; #1},
}

% Attention / Piège (orange)
\newtcolorbox{attentionbox}[1][Attention]{
    ipssibase=ipssiOrange,
    title={\icoAttention\; #1},
}

% Intuition (violet)
\newtcolorbox{intuitionbox}[1][Intuition]{
    ipssibase=ipssiViolet,
    title={\icoIntuition\; #1},
}

% Exemple (cyan)
\newtcolorbox{exemplebox}[1][Exemple]{
    ipssibase=ipssiCyan,
    title={\icoExemple\; #1},
}

% Code Python (style terminal)
\newtcolorbox{pythonbox}[1][Code Python]{
    enhanced,
    breakable,
    left=4mm, right=4mm, top=3mm, bottom=3mm,
    boxrule=0.5pt,
    colback=ipssiGrisClair,
    colframe=ipssiGrisFonce!40,
    fonttitle=\bfseries\ttfamily\small,
    coltitle=ipssiGrisFonce,
    title={\icoPython\; #1},
    attach boxed title to top left={yshift=-2mm, xshift=4mm},
    boxed title style={colback=ipssiGrisClair, colframe=ipssiGrisFonce!40, boxrule=0.5pt, arc=2pt},
    arc=2pt, outer arc=2pt,
}

% ============================================================================
% CODE PYTHON (listings)
% ============================================================================
\usepackage{listings}

\lstdefinestyle{python}{
    language=Python,
    basicstyle=\ttfamily\small,
    keywordstyle=\color{ipssiBleu}\bfseries,
    stringstyle=\color{ipssiVert!70!black},
    commentstyle=\color{gray}\itshape,
    numberstyle=\tiny\color{gray},
    numbers=left,
    numbersep=8pt,
    backgroundcolor=\color{ipssiGrisClair},
    frame=none,
    rulecolor=\color{ipssiGrisFonce!20},
    breaklines=true,
    breakatwhitespace=true,
    tabsize=4,
    showstringspaces=false,
    captionpos=b,
    aboveskip=8pt,
    belowskip=8pt,
    xleftmargin=10pt,
    framexleftmargin=10pt,
    morekeywords={as, with, True, False, None, self, np, plt, import, from},
    literate=
        {é}{{\'e}}1 {è}{{\`e}}1 {ê}{{\^e}}1 {ë}{{\"e}}1
        {à}{{\`a}}1 {â}{{\^a}}1
        {ù}{{\`u}}1 {û}{{\^u}}1
        {ô}{{\^o}}1
        {î}{{\^i}}1 {ï}{{\"i}}1
        {ç}{{\c{c}}}1
        {É}{{\'E}}1 {È}{{\`E}}1
        {→}{{$\rightarrow$}}1
        {≈}{{$\approx$}}1
        {≤}{{$\leq$}}1
        {≥}{{$\geq$}}1
        {★}{{$\bigstar$}}1
        {ŷ}{{$\hat{y}$}}1
        {·}{{$\cdot$}}1,
}

\lstset{style=python}

% Style pour le code inline
\newcommand{\pyinline}[1]{\lstinline[style=python,basicstyle=\ttfamily\small]{#1}}

% ============================================================================
% GRAPHIQUES
% ============================================================================
\usepackage{tikz}
\usepackage{pgfplots}
\pgfplotsset{compat=1.18}
\usetikzlibrary{arrows.meta, calc, positioning, shapes.geometric, decorations.pathreplacing}

% ============================================================================
% TABLEAUX
% ============================================================================
\usepackage{booktabs}
\usepackage{tabularx}
\usepackage{multirow}

% ============================================================================
% LISTES
% ============================================================================
\usepackage{enumitem}
\setlist[itemize]{itemsep=2pt, parsep=0pt, topsep=4pt}
\setlist[enumerate]{itemsep=2pt, parsep=0pt, topsep=4pt}

% ============================================================================
% HYPERLIENS
% ============================================================================
\usepackage[
    colorlinks=true,
    linkcolor=ipssiBleu!80!black,
    urlcolor=ipssiBleu!80!black,
    citecolor=ipssiVert!80!black,
    bookmarks=true,
    bookmarksopen=true,
    pdfauthor={Mohamed EL AFRIT},
    pdftitle={Mathématiques appliquées à l'IA — IPSSI},
    pdfsubject={Formation Bachelor 2 — IPSSI 2025-2026},
]{hyperref}

% ============================================================================
% IMAGES
% ============================================================================
\usepackage{graphicx}

% ============================================================================
% VARIABLES PARAMÉTRABLES (jour et thème)
% ============================================================================
\newcommand{\leJour}{X}
\newcommand{\leTheme}{Thème}
\newcommand{\jourNumero}[1]{\renewcommand{\leJour}{#1}}
\newcommand{\jourTheme}[1]{\renewcommand{\leTheme}{#1}}

% ============================================================================
% EN-TÊTE ET PIED DE PAGE (fancyhdr)
% ============================================================================
\usepackage{fancyhdr}

\pagestyle{fancy}
\fancyhf{}

% En-tête
\fancyhead[L]{\small\sffamily\textcolor{gray}{IPSSI — Maths appliquées à l'IA}}
\fancyhead[R]{\small\sffamily\textcolor{gray}{Jour \leJour\ — \leTheme}}
\renewcommand{\headrulewidth}{0.4pt}

% Pied de page
\fancyfoot[L]{\footnotesize\textcolor{gray}{Mohamed EL AFRIT\enspace\textbullet\enspace\url{www.mohamedelafrit.com}\enspace\textbullet\enspace CC BY-NC-SA 4.0}}
\fancyfoot[C]{\footnotesize\textcolor{gray}{\thepage}}
\fancyfoot[R]{\footnotesize\textcolor{gray}{2025–2026}}
\renewcommand{\footrulewidth}{0.2pt}

% Style pour la page de garde (pas d'en-tête)
\fancypagestyle{pagegarde}{
    \fancyhf{}
    \renewcommand{\headrulewidth}{0pt}
    \fancyfoot[C]{\footnotesize\textcolor{gray}{\thepage}}
    \renewcommand{\footrulewidth}{0pt}
}

% ============================================================================
% PAGE DE GARDE — \pagegarde{titre}{jour}{theme}
% ============================================================================
\newcommand{\pagegarde}[3]{%
    \thispagestyle{pagegarde}
    \begin{center}
        \vspace*{1.5cm}

        % Logo / Bandeau institution
        {\Large\sffamily\bfseries\textcolor{ipssiBleu}{IPSSI}}\par
        \vspace{2pt}
        {\normalsize\sffamily 2\textsuperscript{e} année Bachelor Informatique}\par

        \vspace{0.8cm}
        {\color{ipssiBleu}\rule{0.6\textwidth}{1.5pt}}\par
        \vspace{0.8cm}

        % Titre du module
        {\LARGE\sffamily\bfseries Mathématiques appliquées\\[4pt] à l'Intelligence Artificielle}\par

        \vspace{1cm}

        % Titre du document
        \begin{tcolorbox}[
            enhanced,
            colback=ipssiBleu!5,
            colframe=ipssiBleu!50,
            boxrule=1pt,
            arc=4pt,
            width=0.85\textwidth,
            halign=center,
            top=10pt, bottom=10pt,
        ]
            {\huge\sffamily\bfseries\textcolor{ipssiBleu!80!black}{#1}}\par
            \vspace{6pt}
            {\Large\sffamily Jour #2 — #3}
        \end{tcolorbox}

        \vspace{1.5cm}

        % Informations auteur
        {\large\sffamily
            \textbf{Auteur :} Mohamed EL AFRIT\\[4pt]
            \url{www.mohamedelafrit.com}\\[8pt]
            \textbf{Année académique :} 2025–2026
        }\par

        \vspace{1.5cm}

        % Licence Creative Commons
        {\color{ipssiBleu}\rule{0.4\textwidth}{0.5pt}}\par
        \vspace{8pt}
        {\small\sffamily
            \textcopyright\ 2025 Mohamed EL AFRIT — IPSSI\\[2pt]
            Ce contenu est distribué sous licence\\[2pt]
            \textbf{Creative Commons BY-NC-SA 4.0}\\[2pt]
            \url{https://creativecommons.org/licenses/by-nc-sa/4.0/deed.fr}
        }\par
    \end{center}
    \newpage
}

% ============================================================================
% NIVEAUX D'EXERCICES
% ============================================================================

% Étoile compatible pdflatex
\newcommand{\starfull}[1]{\textcolor{#1}{$\bigstar$}}
\newcommand{\starempty}{\textcolor{gray!40}{$\bigstar$}}

% Fondamental (1 étoile)
\newcommand{\exofondamental}{%
    \starfull{ipssiVert}\starempty\starempty\enspace%
    {\small\sffamily\textcolor{ipssiVert}{\textbf{Fondamental}}}%
}

% Intermédiaire (2 étoiles)
\newcommand{\exointermediaire}{%
    \starfull{ipssiOrange}\starfull{ipssiOrange}\starempty\enspace%
    {\small\sffamily\textcolor{ipssiOrange}{\textbf{Interm\'{e}diaire}}}%
}

% Avancé (3 étoiles)
\newcommand{\exoavance}{%
    \starfull{ipssiRouge}\starfull{ipssiRouge}\starfull{ipssiRouge}\enspace%
    {\small\sffamily\textcolor{ipssiRouge}{\textbf{Avanc\'{e}}}}%
}

% Commande d'exercice numéroté
\newcounter{exocount}[section]
\newcommand{\exercice}[1]{%
    \stepcounter{exocount}%
    \vspace{8pt}%
    \noindent\textbf{\sffamily Exercice \theexocount}\enspace — \enspace #1\par
    \vspace{4pt}%
}

% ============================================================================
% COMMANDES UTILITAIRES
% ============================================================================

% Séparateur visuel entre sections
\newcommand{\separateur}{%
    \vspace{6pt}
    \begin{center}
        {\color{ipssiBleu!30}\rule{0.3\textwidth}{0.5pt}}
    \end{center}
    \vspace{6pt}
}

% Mot-clé en gras coloré
\newcommand{\motcle}[1]{\textbf{\textcolor{ipssiBleu!80!black}{#1}}}

% Formule encadrée importante
\newcommand{\formule}[1]{%
    \begin{center}
        \fcolorbox{ipssiBleu!50}{ipssiBleu!5}{%
            \quad$\displaystyle #1$\quad%
        }
    \end{center}
}

% ============================================================================
% FIN DU PRÉAMBULE
% ============================================================================


\jourNumero{4}
\jourTheme{R\'{e}seaux de neurones}

\begin{document}

\pagegarde{Corrections d\'{e}taill\'{e}es}{4}{R\'{e}seaux de neurones}

\setstretch{1.06}
\setlength{\parskip}{2pt}
\setlength{\abovedisplayskip}{5pt}
\setlength{\belowdisplayskip}{5pt}

\begin{tcolorbox}[
    enhanced, colback=ipssiCyan!5, colframe=ipssiCyan!50,
    boxrule=0.8pt, arc=4pt,
    title={\sffamily\bfseries Rappels},
    coltitle=white, colbacktitle=ipssiCyan!70,
]
\small
$\sigma(z) = \frac{1}{1+e^{-z}}$ \;\;
$\text{ReLU}(z) = \max(0,z)$ \;\;
$\sigma'(z) = \sigma(z)(1-\sigma(z))$ \;\;
$\sigma(0) = 0{,}5$ \;\;
$\sigma(1) \approx 0{,}73$ \;\;
$\sigma(2) \approx 0{,}88$ \;\;
$\sigma(-1) \approx 0{,}27$
\end{tcolorbox}

% ======================================================================
% EXERCICE 1
% ======================================================================
\section*{Exercice 1 \hfill \exofondamental}
\addcontentsline{toc}{section}{Exercice 1 --- Sortie d'un neurone}
\textit{$\vw = (0{,}3;\; 0{,}5)^T$, $b = -2$, activation sigmo\"{\i}de.}

\separateur

\textbf{a)} Somme pond\'{e}r\'{e}e pour $\vx = (5, 4)^T$ :
\[
    z = 0{,}3 \times 5 + 0{,}5 \times 4 + (-2) = 1{,}5 + 2{,}0 - 2{,}0 = \boxed{1{,}5}
\]

\textbf{b)} Activation :
$a = \sigma(1{,}5) = \frac{1}{1 + e^{-1{,}5}} = \frac{1}{1 + 0{,}2231} = \frac{1}{1{,}2231} = \boxed{0{,}8176}$

\textbf{c)} $a = 0{,}818 > 0{,}5$ donc le neurone pr\'{e}dit la \textbf{classe 1}.

\textbf{d)} Pour $\vx = (1, 2)^T$ :
$z = 0{,}3 \times 1 + 0{,}5 \times 2 - 2 = 0{,}3 + 1{,}0 - 2{,}0 = -0{,}7$.

$a = \sigma(-0{,}7) = \frac{1}{1 + e^{0{,}7}} = \frac{1}{1 + 2{,}014} = \boxed{0{,}332}$.
Comme $0{,}332 < 0{,}5$, le neurone pr\'{e}dit la \textbf{classe 0}.

\begin{attentionbox}[Erreurs fr\'{e}quentes]
\begin{itemize}
    \item Oublier le biais $b = -2$ dans le calcul de $z$.
    \item Confondre $e^{-z}$ et $e^{z}$ dans la formule de la sigmo\"{\i}de : c'est $e^{-z}$ au d\'{e}nominateur.
\end{itemize}
\end{attentionbox}

% ======================================================================
% EXERCICE 2
% ======================================================================
\section*{Exercice 2 \hfill \exofondamental}
\addcontentsline{toc}{section}{Exercice 2 --- ReLU et sigmo\"{\i}de}

\separateur

\textbf{a) Sigmo\"{\i}de :}

$\sigma(0) = \frac{1}{1+e^0} = \frac{1}{2} = \boxed{0{,}5}$

$\sigma(2) = \frac{1}{1+e^{-2}} = \frac{1}{1+0{,}135} = \boxed{0{,}881}$

$\sigma(-3) = \frac{1}{1+e^{3}} = \frac{1}{1+20{,}09} = \boxed{0{,}047}$

\textbf{b) ReLU :}

$\text{ReLU}(2{,}5) = \max(0,\; 2{,}5) = \boxed{2{,}5}$

$\text{ReLU}(0) = \max(0,\; 0) = \boxed{0}$

$\text{ReLU}(-1{,}7) = \max(0,\; -1{,}7) = \boxed{0}$

\textbf{c)} $\text{ReLU}(-1) = \max(0, -1) = 0$, puis $\sigma(0) = 0{,}5$.
Donc $\sigma(\text{ReLU}(-1)) = \boxed{0{,}5}$.

\textbf{d)} $\text{ReLU}(3) = 3$, puis $\sigma(3) = \frac{1}{1+e^{-3}} \approx 0{,}953$.
Donc $\sigma(\text{ReLU}(3)) = \boxed{0{,}953}$.

\textbf{e)} $\sigma(z) = \frac{1}{1+e^{-z}}$. Le d\'{e}nominateur $1 + e^{-z}$ est toujours $> 1$
(car $e^{-z} > 0$ pour tout $z$), donc $\sigma(z) < 1$.
De m\^{e}me, $\sigma(z) > 0$ car le num\'{e}rateur et le d\'{e}nominateur sont positifs.
La sigmo\"{\i}de est \textbf{asymptotique} : elle s'approche de 0 et 1 sans jamais les atteindre.

\begin{retenirbox}
    ReLU ``coupe'' les valeurs n\'{e}gatives \`{a} z\'{e}ro. La sigmo\"{\i}de ``\'{e}crase'' tout nombre r\'{e}el dans $(0, 1)$.
    Composer les deux cr\'{e}e une fonction non lin\'{e}aire.
\end{retenirbox}

% ======================================================================
% EXERCICE 3
% ======================================================================
\section*{Exercice 3 \hfill \exofondamental}
\addcontentsline{toc}{section}{Exercice 3 --- Forward pass 2-2-1}
\textit{$\vx = (3, 5)^T$, $\mathbf{W}^{(1)}$, $\mathbf{w}^{(2)}$ donn\'{e}s.}

\separateur

\textbf{a) Sommes pond\'{e}r\'{e}es couche cach\'{e}e :}
\begin{align*}
    z_1 &= 0{,}4 \times 3 + 0{,}6 \times 5 + (-1) = 1{,}2 + 3{,}0 - 1 = \boxed{3{,}2} \\
    z_2 &= (-0{,}3) \times 3 + 0{,}5 \times 5 + 0{,}5 = -0{,}9 + 2{,}5 + 0{,}5 = \boxed{2{,}1}
\end{align*}

\textbf{b) Activations :}
$h_1 = \sigma(3{,}2) \approx \boxed{0{,}961}$ \quad;\quad $h_2 = \sigma(2{,}1) \approx \boxed{0{,}891}$

\textbf{c) Somme pond\'{e}r\'{e}e sortie :}
\[
    z_{\text{out}} = 0{,}7 \times 0{,}961 + 0{,}8 \times 0{,}891 + (-0{,}5) = 0{,}673 + 0{,}713 - 0{,}5 = \boxed{0{,}886}
\]

\textbf{d) Pr\'{e}diction :}
$\yhat = \sigma(0{,}886) \approx \boxed{0{,}708}$

\textbf{e)} $0{,}708 > 0{,}5$ donc le r\'{e}seau pr\'{e}dit la \textbf{classe 1}
(salaire $>$ 45\,k\texteuro).

\begin{attentionbox}[Erreurs fr\'{e}quentes]
\begin{itemize}
    \item Utiliser $\vx$ directement dans la couche de sortie au lieu de $\mathbf{h}$ --- la sortie utilise les \textbf{activations de la couche cach\'{e}e}, pas les entr\'{e}es.
    \item Oublier d'appliquer $\sigma$ \`{a} chaque couche.
\end{itemize}
\end{attentionbox}

% ======================================================================
% EXERCICE 4
% ======================================================================
\section*{Exercice 4 \hfill \exofondamental}
\addcontentsline{toc}{section}{Exercice 4 --- Dimensions des matrices}
\textit{R\'{e}seau 5 $\to$ 8 $\to$ 3 $\to$ 1.}

\separateur

\textbf{a)} $\mathbf{W}^{(1)} \in \mathbb{R}^{8 \times 5}$ (8 neurones, 5 entr\'{e}es).
$\mathbf{b}^{(1)} \in \mathbb{R}^{8}$.

\textbf{b)} $\mathbf{W}^{(2)} \in \mathbb{R}^{3 \times 8}$ (3 neurones, 8 entr\'{e}es).
$\mathbf{b}^{(2)} \in \mathbb{R}^{3}$.

\textbf{c)} $\mathbf{w}^{(3)} \in \mathbb{R}^{3}$ (ou $\mathbf{W}^{(3)} \in \mathbb{R}^{1 \times 3}$).
$b^{(3)} \in \mathbb{R}$.

\textbf{d)} Total des param\`{e}tres :
\begin{align*}
    &\underbrace{8 \times 5}_{W^{(1)}} + \underbrace{8}_{b^{(1)}} + \underbrace{3 \times 8}_{W^{(2)}} + \underbrace{3}_{b^{(2)}} + \underbrace{1 \times 3}_{w^{(3)}} + \underbrace{1}_{b^{(3)}} \\
    &= 40 + 8 + 24 + 3 + 3 + 1 = \boxed{79 \text{ param\`{e}tres}}
\end{align*}

\textbf{e)} Couche 1 passe de 8 \`{a} 16 neurones :
$W^{(1)}$ passe de $8 \times 5 = 40$ \`{a} $16 \times 5 = 80$ ($+40$).
$b^{(1)}$ passe de 8 \`{a} 16 ($+8$).
$W^{(2)}$ passe de $3 \times 8 = 24$ \`{a} $3 \times 16 = 48$ ($+24$).
Total suppl\'{e}mentaire : $\boxed{+72}$ param\`{e}tres (de 79 \`{a} 151).

\begin{retenirbox}
    Le nombre de param\`{e}tres cro\^{i}t \textbf{quadratiquement} avec la taille des couches.
    C'est pourquoi les grands r\'{e}seaux (GPT, etc.)\ ont des milliards de param\`{e}tres.
\end{retenirbox}

% ======================================================================
% EXERCICE 5
% ======================================================================
\section*{Exercice 5 \hfill \exointermediaire}
\addcontentsline{toc}{section}{Exercice 5 --- Composer deux fonctions}
\textit{$f_A(x) = \text{ReLU}(2x-1)$, $g_B(x) = \sigma(3x-2)$.}

\separateur

\textbf{a)} $f_A(0) = \text{ReLU}(-1) = \boxed{0}$ \;;\;
$f_A(1) = \text{ReLU}(1) = \boxed{1}$ \;;\;
$f_A(3) = \text{ReLU}(5) = \boxed{5}$.

\textbf{b) Composition $g_B(f_A(x))$ :}

$x = 0$ : $f_A(0) = 0$, puis $g_B(0) = \sigma(3 \times 0 - 2) = \sigma(-2) \approx \boxed{0{,}119}$.

$x = 1$ : $f_A(1) = 1$, puis $g_B(1) = \sigma(3 \times 1 - 2) = \sigma(1) \approx \boxed{0{,}731}$.

$x = 3$ : $f_A(3) = 5$, puis $g_B(5) = \sigma(3 \times 5 - 2) = \sigma(13) \approx \boxed{1{,}000}$.

\textbf{c)} Non, $g_B \circ f_A$ n'est \textbf{pas} lin\'{e}aire.
Elle combine ReLU (lin\'{e}aire par morceaux) et sigmo\"{\i}de (courbe en S).
La composition de fonctions non lin\'{e}aires est non lin\'{e}aire.

\textbf{d) Sans activation :} $f(x) = 2x - 1$, $g(x) = 3x - 2$.
$g(f(x)) = 3(2x - 1) - 2 = 6x - 3 - 2 = \boxed{6x - 5}$.
C'est une fonction \textbf{lin\'{e}aire} --- deux couches sans activation = une seule couche.

\begin{retenirbox}
    C'est la preuve concr\`{e}te : sans activation, empiler des couches ne sert \`{a} rien.
\end{retenirbox}

% ======================================================================
% EXERCICE 6
% ======================================================================
\section*{Exercice 6 \hfill \exointermediaire}
\addcontentsline{toc}{section}{Exercice 6 --- R\'{e}seau sans activation}

\separateur

\textbf{a-b)} Substitution de $\mathbf{h}$ dans $\yhat$ :
\begin{align*}
    \yhat &= \mathbf{w}^{(2)T}\mathbf{h} + b^{(2)} \\
          &= \mathbf{w}^{(2)T}\bigl(\mathbf{W}^{(1)}\vx + \mathbf{b}^{(1)}\bigr) + b^{(2)} \\
          &= \underbrace{\mathbf{w}^{(2)T}\mathbf{W}^{(1)}}_{\tilde{\vw}^T}\vx
             + \underbrace{\mathbf{w}^{(2)T}\mathbf{b}^{(1)} + b^{(2)}}_{\tilde{b}}
\end{align*}
Donc $\yhat = \tilde{\vw}^T \vx + \tilde{b}$ avec $\tilde{\vw} = \mathbf{W}^{(1)T}\mathbf{w}^{(2)}$.

\textbf{c)} C'est un \textbf{mod\`{e}le lin\'{e}aire} --- exactement comme un perceptron simple.
Le r\'{e}seau \`{a} 2 couches sans activation n'est \textbf{pas plus puissant} qu'un seul neurone.

\textbf{d)} Les fonctions d'activation sont indispensables car elles introduisent la
\textbf{non-lin\'{e}arit\'{e}} qui permet de mod\'{e}liser des fronti\`{e}res courbes et de r\'{e}soudre
des probl\`{e}mes non lin\'{e}airement s\'{e}parables (comme le XOR).

% ======================================================================
% EXERCICE 7
% ======================================================================
\section*{Exercice 7 \hfill \exointermediaire}
\addcontentsline{toc}{section}{Exercice 7 --- Surapprentissage}

\separateur

\textbf{a)} Le surapprentissage commence vers l'\textbf{\'{e}poque 200}.
Signe : l'accuracy d'entra\^{i}nement continue d'augmenter (95\% $\to$ 100\%)
tandis que l'accuracy de validation \textbf{diminue} (87\% $\to$ 75\%).
L'\'{e}cart train/validation se creuse.

\textbf{b)} Le meilleur mod\`{e}le est \`{a} l'\textbf{\'{e}poque 50} : accuracy de validation
maximale (83\%), avec un \'{e}cart raisonnable par rapport \`{a} l'entra\^{i}nement (85\% vs 83\%).

\textbf{c)} Trois solutions :
\begin{enumerate}
    \item \textbf{Early stopping} : arr\^{e}ter l'entra\^{i}nement quand la validation cesse de s'am\'{e}liorer.
    \item \textbf{R\'{e}gularisation $L_2$} : ajouter un terme $\lambda\|\vw\|^2$ \`{a} la loss pour p\'{e}naliser les poids trop grands.
    \item \textbf{Dropout} : d\'{e}sactiver al\'{e}atoirement des neurones pendant l'entra\^{i}nement pour forcer la redondance.
\end{enumerate}

\textbf{d)} 100\% sur l'entra\^{i}nement signifie que le mod\`{e}le a \textbf{m\'{e}moris\'{e}} toutes les observations, y compris le bruit et les cas atypiques. C'est comme un \'{e}l\`{e}ve qui apprend ses fiches par c\oe ur sans comprendre : il \'{e}choue face \`{a} un probl\`{e}me nouveau.

\begin{attentionbox}[Erreurs fr\'{e}quentes]
\begin{itemize}
    \item Choisir le mod\`{e}le avec la meilleure accuracy d'\textbf{entra\^{i}nement} --- il faut regarder la \textbf{validation}.
    \item Croire qu'un mod\`{e}le plus complexe est toujours meilleur : au-del\`{a} d'un certain point, il surapprentit.
\end{itemize}
\end{attentionbox}

% ======================================================================
% EXERCICE 8
% ======================================================================
\section*{Exercice 8 \hfill \exointermediaire}
\addcontentsline{toc}{section}{Exercice 8 --- Sigmo\"{\i}de et ReLU (Python)}

\separateur

\begin{pythonbox}[Correction compl\`{e}te]
\begin{lstlisting}
import numpy as np
import matplotlib.pyplot as plt

# 1. Sigmoide
def sigmoid(z):
    return 1 / (1 + np.exp(-z))

# 2. ReLU
def relu(z):
    return np.maximum(0, z)

# 3. Test
valeurs = np.array([-3, -1, 0, 1, 2, 5])
print("z       :", valeurs)
print("sigmoid :", np.round(sigmoid(valeurs), 4))
print("relu    :", relu(valeurs))

# 4. Courbes
z = np.linspace(-6, 6, 200)
fig, (ax1, ax2) = plt.subplots(1, 2, figsize=(12, 4))

ax1.plot(z, sigmoid(z), 'b-', lw=2)
ax1.set_title("Sigmoide"); ax1.grid(True, alpha=0.3)
ax1.axhline(0.5, color='gray', ls=':')

ax2.plot(z, relu(z), 'r-', lw=2)
ax2.set_title("ReLU"); ax2.grid(True, alpha=0.3)

plt.tight_layout(); plt.show()

# 5. Derivee
def sigmoid_deriv(z):
    s = sigmoid(z)
    return s * (1 - s)

print("\nDerivee sigmoid(0) =", sigmoid_deriv(0))
print("Derivee sigmoid(2) =", round(sigmoid_deriv(2), 4))
\end{lstlisting}
\end{pythonbox}

\textbf{Sortie attendue :}
\begin{verbatim}
z       : [-3 -1  0  1  2  5]
sigmoid : [0.0474 0.2689 0.5    0.7311 0.8808 0.9933]
relu    : [0 0 0 1 2 5]

Derivee sigmoid(0) = 0.25
Derivee sigmoid(2) = 0.1050
\end{verbatim}

\begin{intuitionbox}[Variante]
    Essayez aussi \texttt{np.tanh(z)} (tangente hyperbolique, sortie dans $(-1, 1)$)
    et la \textbf{Leaky ReLU} : \texttt{np.where(z > 0, z, 0.01 * z)}.
\end{intuitionbox}

% ======================================================================
% EXERCICE 9
% ======================================================================
\section*{Exercice 9 \hfill \exointermediaire}
\addcontentsline{toc}{section}{Exercice 9 --- Forward pass 2-2-1 (Python)}

\separateur

\begin{pythonbox}[Correction compl\`{e}te]
\begin{lstlisting}
import numpy as np

def sigmoid(z):
    return 1 / (1 + np.exp(-z))

x = np.array([5.0, 4.0])

W1 = np.array([[0.3, 0.5], [-0.2, 0.4]])
b1 = np.array([-2.0, 0.5])
w2 = np.array([0.8, 0.6])
b2 = -0.5

# 1. Couche cachee
z1 = W1 @ x + b1
h  = sigmoid(z1)
print(f"z1 = {z1}")
print(f"h  = {np.round(h, 4)}")

# 2. Couche de sortie
z2 = w2 @ h + b2
y_pred = sigmoid(z2)
print(f"\nz2 = {z2:.4f}")
print(f"y_pred = {y_pred:.4f}")

# 3. Decision
classe = 1 if y_pred > 0.5 else 0
print(f"\nClasse predite : {classe}")

# 4. Test sur 6 developpeurs
print("\n--- Predictions ---")
devs = np.array([[0,2],[1,1],[5,4],[9,5],[12,7],[18,5]])
for d in devs:
    h_i = sigmoid(W1 @ d + b1)
    y_i = sigmoid(w2 @ h_i + b2)
    print(f"  x={d} -> y={y_i:.3f} -> "
          f"classe {1 if y_i > 0.5 else 0}")
\end{lstlisting}
\end{pythonbox}

\textbf{Sortie attendue :}
\begin{verbatim}
z1 = [1.5 1.1]
h  = [0.8176 0.7503]

z2 = 0.6043
y_pred = 0.6466

Classe predite : 1

--- Predictions ---
  x=[0 2] -> y=0.551 -> classe 1
  x=[1 1] -> y=0.533 -> classe 1
  x=[5 4] -> y=0.647 -> classe 1
  x=[9 5] -> y=0.693 -> classe 1
  x=[12 7] -> y=0.726 -> classe 1
  x=[18 5] -> y=0.715 -> classe 1
\end{verbatim}

\begin{attentionbox}[Erreurs fr\'{e}quentes]
\begin{itemize}
    \item \'{E}crire \texttt{w2 @ x + b2} au lieu de \texttt{w2 @ h + b2} --- la sortie utilise \texttt{h}, pas \texttt{x}.
    \item Avec ces poids non entra\^{i}n\'{e}s, le r\'{e}seau pr\'{e}dit ``classe~1'' pour tout le monde --- c'est normal, il n'a pas appris.
\end{itemize}
\end{attentionbox}

% ======================================================================
% EXERCICE 10
% ======================================================================
\section*{Exercice 10 \hfill \exointermediaire}
\addcontentsline{toc}{section}{Exercice 10 --- XOR avec backpropagation (Python)}

\separateur

\begin{pythonbox}[Correction compl\`{e}te]
\begin{lstlisting}
import numpy as np

def sigmoid(z):
    return 1 / (1 + np.exp(-z))
def sigmoid_d(a):
    return a * (1 - a)

X = np.array([[0,0],[0,1],[1,0],[1,1]], dtype=float)
y = np.array([[0],[1],[1],[0]], dtype=float)

np.random.seed(42)
W1 = np.random.randn(2, 4) * 0.5
b1 = np.zeros((1, 4))
W2 = np.random.randn(4, 1) * 0.5
b2 = np.zeros((1, 1))

alpha = 1.0

for epoch in range(5000):
    # Forward
    z1 = X @ W1 + b1
    h  = sigmoid(z1)
    z2 = h @ W2 + b2
    out = sigmoid(z2)

    loss = np.mean((y - out) ** 2)

    # Backward
    d_out = (out - y) * sigmoid_d(out)
    dW2 = h.T @ d_out / 4
    db2 = np.mean(d_out, axis=0, keepdims=True)

    d_h = d_out @ W2.T * sigmoid_d(h)
    dW1 = X.T @ d_h / 4
    db1 = np.mean(d_h, axis=0, keepdims=True)

    # Mise a jour
    W2 -= alpha * dW2; b2 -= alpha * db2
    W1 -= alpha * dW1; b1 -= alpha * db1

    if epoch % 1000 == 999:
        preds = (out > 0.5).astype(int).flatten()
        acc = np.mean(preds == y.flatten())
        print(f"Epoch {epoch+1} : "
              f"loss={loss:.4f}, acc={acc:.0%}")

print(f"\nPredictions : {out.flatten().round(3)}")
print(f"Attendu     : {y.flatten().astype(int)}")
\end{lstlisting}
\end{pythonbox}

\textbf{Sortie attendue :}
\begin{verbatim}
Epoch 1000 : loss=0.2336, acc=75%
Epoch 2000 : loss=0.0124, acc=100%
Epoch 3000 : loss=0.0048, acc=100%
Epoch 4000 : loss=0.0027, acc=100%
Epoch 5000 : loss=0.0018, acc=100%

Predictions : [0.039 0.965 0.965 0.038]
Attendu     : [0 1 1 0]
\end{verbatim}

\begin{retenirbox}
    Le r\'{e}seau 2-4-1 r\'{e}sout le XOR avec 100\% d'accuracy !
    Les 4 neurones cach\'{e}s cr\'{e}ent des \textbf{repr\'{e}sentations interm\'{e}diaires}
    qui rendent le probl\`{e}me lin\'{e}airement s\'{e}parable dans $\mathbb{R}^4$.
\end{retenirbox}

% ======================================================================
% EXERCICE 11
% ======================================================================
\section*{Exercice 11 \hfill \exoavance}
\addcontentsline{toc}{section}{Exercice 11 --- R\'{e}seau from scratch salaire (Python)}

\separateur

\begin{pythonbox}[Correction compl\`{e}te]
\begin{lstlisting}
import numpy as np
import matplotlib.pyplot as plt

def sigmoid(z):
    return 1 / (1 + np.exp(-np.clip(z, -500, 500)))
def sigmoid_d(a):
    return a * (1 - a)

# a) Dataset
data = np.array([
    [0,2,28],[1,1,26.5],[1,3,35],[2,2,32],[2,4,38],
    [0,1,25],[3,3,34.5],[1,5,42],[4,3,38.5],[5,4,44],
    [5,3,41],[6,5,48],[6,4,43.5],[7,5,50],[7,3,39],
    [4,6,47],[8,4,51],[5,2,40],[9,5,52],[10,6,55],
    [10,4,48.5],[11,5,58],[12,7,60],[9,3,42],
    [13,6,62],[11,4,54],[15,8,65],[17,6,68],
    [20,7,72],[18,5,58]], dtype=float)

X_raw = data[:, :2]
y = (data[:, 2] > 45).astype(float).reshape(-1, 1)

# b) Normalisation
mu = X_raw.mean(axis=0)
sigma_x = X_raw.std(axis=0)
X = (X_raw - mu) / sigma_x

# c) Reseau 2-4-1
np.random.seed(42)
W1 = np.random.randn(2, 4) * 0.5
b1 = np.zeros((1, 4))
W2 = np.random.randn(4, 1) * 0.5
b2 = np.zeros((1, 1))
alpha = 0.5
n = len(y)
hist_loss = []

for epoch in range(2000):
    h = sigmoid(X @ W1 + b1)
    out = sigmoid(h @ W2 + b2)
    loss = np.mean((y - out) ** 2)
    hist_loss.append(loss)

    d_out = (out - y) * sigmoid_d(out)
    W2 -= alpha * (h.T @ d_out) / n
    b2 -= alpha * np.mean(d_out, axis=0, keepdims=True)
    d_h = (d_out @ W2.T) * sigmoid_d(h)
    W1 -= alpha * (X.T @ d_h) / n
    b1 -= alpha * np.mean(d_h, axis=0, keepdims=True)

    if epoch % 500 == 499:
        acc = np.mean((out > 0.5) == y)
        print(f"Epoch {epoch+1:5d} : "
              f"loss={loss:.4f}, acc={acc:.0%}")

# e) Graphiques
fig, (ax1, ax2) = plt.subplots(1, 2, figsize=(14, 5))
ax1.plot(hist_loss, 'b-', lw=1)
ax1.set_xlabel("Epoch"); ax1.set_ylabel("Loss")
ax1.set_title("Convergence"); ax1.grid(True, alpha=0.3)

# Surface de decision
xx, yy = np.meshgrid(
    np.linspace(X[:,0].min()-1, X[:,0].max()+1, 100),
    np.linspace(X[:,1].min()-1, X[:,1].max()+1, 100))
grid = np.c_[xx.ravel(), yy.ravel()]
h_g = sigmoid(grid @ W1 + b1)
z_g = sigmoid(h_g @ W2 + b2).reshape(xx.shape)
ax2.contourf(xx, yy, z_g, levels=20,
             cmap='RdYlGn', alpha=0.5)
ax2.contour(xx, yy, z_g, levels=[0.5],
            colors='black', linewidths=2)
for c, col in [(0,'#e74c3c'),(1,'#2ecc71')]:
    mask = y.flatten() == c
    ax2.scatter(X[mask,0], X[mask,1], c=col,
                s=60, edgecolors='white', lw=0.5)
ax2.set_title("Surface de decision")
ax2.set_xlabel("Exp (norm)")
ax2.set_ylabel("Lang (norm)")
plt.tight_layout(); plt.show()
\end{lstlisting}
\end{pythonbox}

\textbf{Sortie attendue :}
\begin{verbatim}
Epoch   500 : loss=0.1073, acc=90%
Epoch  1000 : loss=0.0681, acc=93%
Epoch  1500 : loss=0.0523, acc=93%
Epoch  2000 : loss=0.0435, acc=97%
\end{verbatim}

L'accuracy (93--97\%) d\'{e}passe celle du perceptron (${\sim}87$--$90\%$). La fronti\`{e}re est \textbf{courbe}, ce qui permet de mieux s\'{e}parer les cas atypiques.

\begin{attentionbox}[Erreurs fr\'{e}quentes]
\begin{itemize}
    \item Oublier de \textbf{normaliser} les features --- sans normalisation, le r\'{e}seau converge tr\`{e}s lentement ou pas du tout.
    \item Utiliser \texttt{np.clip} dans la sigmo\"{\i}de pour \'{e}viter les overflow num\'{e}riques sur de grandes valeurs.
\end{itemize}
\end{attentionbox}

% ======================================================================
% EXERCICE 12
% ======================================================================
\section*{Exercice 12 \hfill \exoavance}
\addcontentsline{toc}{section}{Exercice 12 --- NumPy vs MLPClassifier (Python)}

\separateur

\begin{pythonbox}[Correction compl\`{e}te]
\begin{lstlisting}
import numpy as np
import matplotlib.pyplot as plt
from sklearn.neural_network import MLPClassifier
from sklearn.linear_model import Perceptron as SkPerc
from sklearn.preprocessing import StandardScaler

data = np.array([
    [0,2,28],[1,1,26.5],[1,3,35],[2,2,32],[2,4,38],
    [0,1,25],[3,3,34.5],[1,5,42],[4,3,38.5],[5,4,44],
    [5,3,41],[6,5,48],[6,4,43.5],[7,5,50],[7,3,39],
    [4,6,47],[8,4,51],[5,2,40],[9,5,52],[10,6,55],
    [10,4,48.5],[11,5,58],[12,7,60],[9,3,42],
    [13,6,62],[11,4,54],[15,8,65],[17,6,68],
    [20,7,72],[18,5,58]], dtype=float)

X_raw = data[:,:2]
y = (data[:,2] > 45).astype(int)
scaler = StandardScaler()
X = scaler.fit_transform(X_raw)

# M1: From scratch (meme code que ex 11)
def sig(z): return 1/(1+np.exp(-np.clip(z,-500,500)))
def sig_d(a): return a*(1-a)
np.random.seed(42)
W1=np.random.randn(2,4)*0.5; b1=np.zeros((1,4))
W2=np.random.randn(4,1)*0.5; b2=np.zeros((1,1))
yr = y.reshape(-1,1).astype(float)
for _ in range(2000):
    h=sig(X@W1+b1); out=sig(h@W2+b2)
    do=(out-yr)*sig_d(out)
    W2-=0.5*(h.T@do)/30; b2-=0.5*np.mean(do,axis=0,keepdims=True)
    dh=(do@W2.T)*sig_d(h)
    W1-=0.5*(X.T@dh)/30; b1-=0.5*np.mean(dh,axis=0,keepdims=True)
acc_fs = np.mean((out>0.5).flatten()==y)

# M2: MLPClassifier
mlp = MLPClassifier(hidden_layer_sizes=(4,),
    max_iter=2000, activation='logistic',
    random_state=42)
mlp.fit(X, y); acc_mlp = mlp.score(X, y)

# M3: Perceptron
perc = SkPerc(max_iter=100, eta0=0.1,
              random_state=42)
perc.fit(X, y); acc_perc = perc.score(X, y)

# c) Tableau
print(f"{'Methode':<22}|{'Acc':>6}|{'Params':>7}")
print("-" * 38)
print(f"{'From scratch 2-4-1':<22}|{acc_fs:>5.0%}|"
      f"{2*4+4+4*1+1:>7}")
print(f"{'MLPClassifier (4,)':<22}|{acc_mlp:>5.0%}|"
      f"{2*4+4+4*1+1:>7}")
print(f"{'Perceptron':<22}|{acc_perc:>5.0%}|"
      f"{'3':>7}")

# d) 3 surfaces de decision
fig, axes = plt.subplots(1, 3, figsize=(16, 5))
xx, yy = np.meshgrid(
    np.linspace(X[:,0].min()-1, X[:,0].max()+1, 100),
    np.linspace(X[:,1].min()-1, X[:,1].max()+1, 100))
grid = np.c_[xx.ravel(), yy.ravel()]

# From scratch
h_g=sig(grid@W1+b1)
z_fs=sig(h_g@W2+b2).reshape(xx.shape)
# MLPClassifier
z_mlp=mlp.predict_proba(grid)[:,1].reshape(xx.shape)
# Perceptron
z_perc=perc.decision_function(grid).reshape(xx.shape)

titles = [f'From scratch ({acc_fs:.0%})',
          f'MLPClassifier ({acc_mlp:.0%})',
          f'Perceptron ({acc_perc:.0%})']
zz = [z_fs, z_mlp, z_perc]
lvls = [[0.5], [0.5], [0]]

for ax, z, t, lv in zip(axes, zz, titles, lvls):
    ax.contourf(xx, yy, z, levels=20,
                cmap='RdYlGn', alpha=0.4)
    ax.contour(xx, yy, z, levels=lv,
               colors='black', linewidths=2)
    for c, col in [(0,'#e74c3c'),(1,'#2ecc71')]:
        mask = y == c
        ax.scatter(X[mask,0], X[mask,1], c=col,
                   s=50, edgecolors='white', lw=0.5)
    ax.set_title(t); ax.grid(True, alpha=0.2)
plt.tight_layout(); plt.show()

# e) Le reseau est meilleur que le perceptron.
# Sa frontiere est COURBE (non lineaire).
# Le perceptron a 3 params, le reseau 17.
\end{lstlisting}
\end{pythonbox}

\textbf{Sortie attendue :}
\begin{verbatim}
Methode               |   Acc| Params
--------------------------------------
From scratch 2-4-1    |  97% |     17
MLPClassifier (4,)    |  93% |     17
Perceptron            |  87% |      3
\end{verbatim}

\begin{retenirbox}[Conclusion g\'{e}n\'{e}rale du Jour 4]
    Le r\'{e}seau multicouche est strictement plus puissant que le perceptron :
    sa fronti\`{e}re \textbf{courbe} capture les relations non lin\'{e}aires.
    Le prix \`{a} payer : plus de param\`{e}tres (17 vs 3), plus de temps d'entra\^{i}nement,
    et un risque de surapprentissage. C'est le \textbf{compromis biais--variance} fondamental du ML.
\end{retenirbox}

\begin{intuitionbox}[Bilan des 4 jours]
    Vous avez construit, \textbf{from scratch}, toute la cha\^{i}ne de l'IA :
    \textbf{J1}~vecteurs et matrices, \textbf{J2}~descente de gradient,
    \textbf{J3}~perceptron, \textbf{J4}~r\'{e}seau de neurones.
    Chaque framework moderne (TensorFlow, PyTorch) impl\'{e}mente exactement ces concepts
    --- en plus rapide et plus grand. Vous avez les fondations pour les comprendre.
\end{intuitionbox}

% ======================================================================
\vfill
\begin{center}
    {\footnotesize\sffamily\textcolor{gray}{%
        \textcopyright\ 2025 Mohamed EL AFRIT --- IPSSI \enspace\textbullet\enspace
        \url{www.mohamedelafrit.com} \enspace\textbullet\enspace
        CC BY-NC-SA 4.0
    }}
\end{center}

\end{document}
